\documentclass{article}
\usepackage{amsmath,amssymb,amsthm}
\usepackage{algorithm,algorithmic}
\usepackage{enumitem}
\usepackage{hyperref}
\usepackage{bm}
\newtheorem{theorem}{Theorem}
\newtheorem{lemma}{Lemma}
\newtheorem{assumption}{Assumption}
\newtheorem{corollary}{Corollary}
\newcommand{\inner}[2]{\langle #1,#2\rangle}
\newcommand{\D}[2]{D_{h}(#1,#2)}
\newcommand{\grad}{\nabla}
\newcommand{\prox}{\operatorname{prox}}
\begin{document}

\section*{Algorithm Name}
\textbf{Mirror Descent with Bregman Prox (Relative Gradient Method)}

\section*{Mathematical Formulation}
Let $f:\mathbb{R}^{n}\to\mathbb{R}$ be a convex differentiable objective and let $h:\mathbb{R}^{n}\to\mathbb{R}$ be a \emph{prox‑function} that is $\sigma$–strongly convex with respect to a norm $\|\cdot\|$.
The associated Bregman divergence is  

\[
\D{x}{y}=h(x)-h(y)-\inner{\grad h(y)}{x-y}.
\]

Given a stepsize sequence $\{\eta_k\}_{k\ge 0}$, the \emph{mirror descent} iterate is defined by  

\[
\boxed{\;
x_{k+1}= \arg\min_{x\in\mathbb{R}^{n}}
\Big\{ \inner{\grad f(x_k)}{x} + \frac{1}{\eta_k}\,\D{x}{x_k}\Big\}
\;}
\tag{MD}
\]

The optimality condition of (MD) yields the \emph{mirror map}  

\[
\grad h(x_{k+1}) = \grad h(x_k)-\eta_k \grad f(x_k). \tag{1}
\]

\section*{Pseudocode}
\begin{algorithm}[H]
\caption{Mirror Descent (Bregman Prox)}
\begin{algorithmic}[1]
\REQUIRE Initial point $x_0\in\mathbb{R}^{n}$, stepsizes $\{\eta_k\}_{k\ge0}>0$, prox‑function $h$.
\FOR{$k=0,1,2,\dots$}
    \STATE Compute gradient $g_k=\grad f(x_k)$.
    \STATE Update $x_{k+1}= \arg\min_{x}\Big\{ \inner{g_k}{x}+ \frac{1}{\eta_k}\D{x}{x_k}\Big\}$.
    \STATE (Equivalently, compute $x_{k+1}$ from (1): $\grad h(x_{k+1})=\grad h(x_k)-\eta_k g_k$.)
\ENDFOR
\ENSURE Final iterate $x_T$ (or averaged iterate $\bar x_T$).
\end{algorithmic}
\end{algorithm}

\section*{Dry Run Example}
We minimise $f(w)=(w-3)^2$ on $\mathbb{R}$.
Choose the Euclidean prox $h(w)=\tfrac12 w^{2}$, hence $\D{u}{v}= \tfrac12 (u-v)^2$ and $\grad h(w)=w$.
Set $w_0=0$, constant stepsize $\eta_k=\eta=0.1$.

\begin{enumerate}[label=Iteration \arabic*:]
\item $k=0$  
\[
g_0=\grad f(w_0)=2(w_0-3)=-6.
\]
Update via (1): $w_1 = w_0-\eta g_0 =0-0.1(-6)=0.6.$  
Objective: $f(w_1)=(0.6-3)^2=5.76$.

\item $k=1$  
\[
g_1=2(w_1-3)=2(0.6-3)=-4.8.
\]
Update: $w_2 = w_1-\eta g_1 =0.6-0.1(-4.8)=1.08.$  
Objective: $f(w_2)=(1.08-3)^2=3.6864$.

\item $k=2$  
\[
g_2=2(w_2-3)=2(1.08-3)=-3.84.
\]
Update: $w_3 = w_2-\eta g_2 =1.08-0.1(-3.84)=1.464.$  
Objective: $f(w_3)=(1.464-3)^2=2.3645\ (\text{approx}).
\]
\end{enumerate}
The iterates move monotonically towards the minimiser $w^{*}=3$ and the objective value decreases.

\newpage
\section*{Assumptions}
\begin{assumption}[Convexity and Relative Smoothness]\label{ass:convex}
The objective $f$ is convex and \emph{L‑smooth relative to $h$}, i.e. for some $L>0$,
\[
f(y)\le f(x)+\inner{\grad f(x)}{y-x}+L\,\D{y}{x},\qquad\forall x,y\in\mathbb{R}^{n}. \tag{2}
\]
\end{assumption}
\begin{assumption}[Strong Convexity of the Prox‑function]\label{ass:strong}
The prox‑function $h$ is $\sigma$‑strongly convex w.r.t. $\|\cdot\|$:
\[
\D{x}{y}\ge\frac{\sigma}{2}\|x-y\|^{2},\qquad\forall x,y\in\mathbb{R}^{n}. \tag{3}
\]
\end{assumption}
\begin{assumption}[Stepsize]\label{ass:step}
The stepsizes satisfy $0<\eta_k\le \frac{1}{L}$ for all $k$.  In the constant‑stepsize case we write $\eta_k=\eta\le1/L$.
\end{assumption}
\begin{assumption}[Existence of a Minimiser]\label{ass:opt}
There exists $x^{*}\in\arg\min f$, and the Bregman distance $\D{x^{*}}{x_0}<\infty$.
\end{assumption}

\section*{Main Convergence Theorem}
\begin{theorem}[Sublinear $O(1/k)$ rate for Mirror Descent]\label{thm:rate}
Assume \ref{ass:convex}--\ref{ass:opt} and use a constant stepsize $\eta\le 1/L$.  
For the iterates $\{x_k\}$ generated by \eqref{MD} we have for every $k\ge1$
\[
f(x_k)-f(x^{*})\;\le\;\frac{\D{x^{*}}{x_0}}{\eta\,k}. \tag{4}
\]
Consequently $f(x_k)\to f(x^{*})$ at the rate $O(1/k)$.
\end{theorem}

\section*{Theoretical Properties}
\begin{itemize}
\item \textbf{Stability.} The update (1) is a contraction in the Bregman geometry when $\eta\le 1/L$; the Bregman distance to the optimum never increases.
\item \textbf{Hyperparameter Sensitivity.} The bound (4) depends linearly on $1/\eta$; choosing $\eta$ close to $1/L$ yields the smallest constant while preserving feasibility.
\item \textbf{Comparison to Euclidean Gradient Descent.} When $h(x)=\tfrac12\|x\|^{2}$, $\D{\cdot}{\cdot}$ reduces to the squared Euclidean norm and Theorem~\ref{thm:rate} recovers the classic $O(1/k)$ rate for gradient descent with stepsize $1/L$.
\item \textbf{Extension to Averaged Iterate.} Defining $\bar x_{k}= \frac{1}{k}\sum_{t=1}^{k}x_t$ yields the same bound (4) by convexity of $f$.
\end{itemize}

\section*{Discussion}
The theorem shows that mirror descent inherits the optimal $O(1/k)$ convergence for convex smooth problems when the smoothness is measured relative to the chosen Bregman geometry.  This flexibility enables the algorithm to exploit problem structure (e.g., simplex constraints via negative entropy) while keeping the same theoretical guarantees as Euclidean gradient descent.

\newpage
\section*{Preliminaries (Definitions and Lemmas)}
\begin{lemma}[Three‑Point Identity]\label{lem:3pt}
For any $x,y,z$,
\[
\inner{\grad h(z)-\grad h(y)}{x-z}= \D{x}{y}-\D{x}{z}-\D{z}{y}. \tag{5}
\]
\end{lemma}
\begin{proof}
Direct expansion of the definition of $\D{\cdot}{\cdot}$ yields (5). \hfill\qedsymbol
\end{proof}

\begin{lemma}[Relative Smoothness Inequality]\label{lem:relSmooth}
Under Assumption~\ref{ass:convex}, for any $x,y$,
\[
f(y)\le f(x)+\inner{\grad f(x)}{y-x}+L\,\D{y}{x}. \tag{6}
\]
\end{lemma}
\begin{proof}
This is precisely the statement of relative smoothness. \hfill\qedsymbol
\end{proof}

\section*{Main Proof (Step‑by‑Step Derivation)}
We prove Theorem~\ref{thm:rate}.  
All non‑trivial steps are numbered \textbf{[Step N]}.

\begin{proof}
Let $x^{*}$ be any optimal point.  
From the optimality condition of \eqref{MD} (or equivalently (1)) we have for all $x$
\[
\inner{\grad f(x_k)}{x_{k+1}-x}+ \frac{1}{\eta_k}\inner{\grad h(x_{k+1})-\grad h(x_k)}{x_{k+1}-x}\le 0. \tag{7}
\]

\noindent\textbf{[Step 1]} Choose $x=x^{*}$ in (7) and rearrange:
\[
\inner{\grad f(x_k)}{x_{k+1}-x^{*}}
\le -\frac{1}{\eta_k}\inner{\grad h(x_{k+1})-\grad h(x_k)}{x_{k+1}-x^{*}}.
\]

\noindent\textbf{[Step 2]} Apply Lemma~\ref{lem:3pt} with $(x,y,z)=(x^{*},x_k,x_{k+1})$:
\[
\inner{\grad h(x_{k+1})-\grad h(x_k)}{x^{*}-x_{k+1}}
= \D{x^{*}}{x_k}-\D{x^{*}}{x_{k+1}}-\D{x_{k+1}}{x_k}. \tag{8}
\]

\noindent\textbf{[Step 3]} Substitute (8) into the inequality of Step 1 and multiply by $1/\eta_k$:
\[
\inner{\grad f(x_k)}{x_{k+1}-x^{*}}
\le \frac{1}{\eta_k}\Big(\D{x^{*}}{x_{k+1}}-\D{x^{*}}{x_k}+\D{x_{k+1}}{x_k}\Big). \tag{9}
\]

\noindent\textbf{[Step 4]} By convexity of $f$,
\[
f(x_{k})-f(x^{*})\le \inner{\grad f(x_k)}{x_{k}-x^{*}}. \tag{10}
\]

\noindent\textbf{[Step 5]} Add and subtract $x_{k+1}$ inside the inner product of (10):
\[
\inner{\grad f(x_k)}{x_{k}-x^{*}}
= \inner{\grad f(x_k)}{x_{k}-x_{k+1}}+\inner{\grad f(x_k)}{x_{k+1}-x^{*}}. \tag{11}
\]

\noindent\textbf{[Step 6]} Combine (9) and (11):
\[
f(x_{k})-f(x^{*})
\le \inner{\grad f(x_k)}{x_{k}-x_{k+1}}
+\frac{1}{\eta_k}\Big(\D{x^{*}}{x_{k+1}}-\D{x^{*}}{x_k}+\D{x_{k+1}}{x_k}\Big). \tag{12}
\]

\noindent\textbf{[Step 7]} Use relative smoothness Lemma~\ref{lem:relSmooth} with $(x,y)=(x_k,x_{k+1})$:
\[
f(x_{k+1})\le f(x_k)+\inner{\grad f(x_k)}{x_{k+1}-x_k}+L\,\D{x_{k+1}}{x_k}. \tag{13}
\]

\noindent\textbf{[Step 8]} Rearrange (13) to obtain a bound on the inner product:
\[
\inner{\grad f(x_k)}{x_k-x_{k+1}}\le f(x_k)-f(x_{k+1})+L\,\D{x_{k+1}}{x_k}. \tag{14}
\]

\noindent\textbf{[Step 9]} Substitute (14) into (12):
\[
\begin{aligned}
f(x_{k})-f(x^{*})
&\le f(x_k)-f(x_{k+1})+L\,\D{x_{k+1}}{x_k}\\
&\quad +\frac{1}{\eta_k}\Big(\D{x^{*}}{x_{k+1}}-\D{x^{*}}{x_k}+\D{x_{k+1}}{x_k}\Big).
\end{aligned} \tag{15}
\]

\noindent\textbf{[Step 10]} Cancel $f(x_k)$ on both sides and collect the Bregman terms:
\[
f(x_{k+1})-f(x^{*})
\le \frac{1}{\eta_k}\big(\D{x^{*}}{x_k}-\D{x^{*}}{x_{k+1}}\big)
+\Big(L-\frac{1}{\eta_k}\Big)\D{x_{k+1}}{x_k}. \tag{16}
\]

\noindent\textbf{[Step 11]} By Assumption~\ref{ass:step}, $\eta_k\le 1/L$ implies $L-\frac{1}{\eta_k}\le 0$.  
Hence the second term on the right‑hand side of (16) is non‑positive and can be dropped:
\[
f(x_{k+1})-f(x^{*})
\le \frac{1}{\eta_k}\big(\D{x^{*}}{x_k}-\D{x^{*}}{x_{k+1}}\big). \tag{17}
\]

\noindent\textbf{[Step 12]} Rearranging (17) gives a recursion for the Bregman distance:
\[
\D{x^{*}}{x_{k+1}}\le \D{x^{*}}{x_k}-\eta_k\big(f(x_{k+1})-f(x^{*})\big). \tag{18}
\]

\noindent\textbf{[Step 13]} Sum (18) from $k=0$ to $k=T-1$ (telescoping the Bregman terms):
\[
\D{x^{*}}{x_T}\le \D{x^{*}}{x_0}-\sum_{k=0}^{T-1}\eta_k\big(f(x_{k+1})-f(x^{*})\big). \tag{19}
\]

\noindent\textbf{[Step 14]} Since $\D{x^{*}}{x_T}\ge 0$, we obtain
\[
\sum_{k=0}^{T-1}\eta_k\big(f(x_{k+1})-f(x^{*})\big)\le \D{x^{*}}{x_0}. \tag{20}
\]

\noindent\textbf{[Step 15]} For constant stepsize $\eta_k=\eta$, divide both sides of (20) by $\eta T$:
\[
\frac{1}{T}\sum_{k=1}^{T}\big(f(x_{k})-f(x^{*})\big)
\le \frac{\D{x^{*}}{x_0}}{\eta T}. \tag{21}
\]

\noindent\textbf{[Step 16]} By convexity of $f$, the average of the function values dominates the function value at the last iterate:
\[
f(x_T)-f(x^{*})\le \frac{1}{T}\sum_{k=1}^{T}\big(f(x_{k})-f(x^{*})\big). \tag{22}
\]

\noindent\textbf{[Step 17]} Combining (21) and (22) yields the desired bound (4):
\[
f(x_T)-f(x^{*})\le \frac{\D{x^{*}}{x_0}}{\eta T}.
\qquad\blacksquare
\]
\end{proof}

\section*{Rate Derivation (With Constants)}
From (4) we read the explicit constant:
\[
\boxed{C:=\frac{\D{x^{*}}{x_0}}{\eta}}.
\]
If $h$ is $\sigma$‑strongly convex (Assumption~\ref{ass:strong}) and the norm is Euclidean, then
\[
\D{x^{*}}{x_0}\ge \frac{\sigma}{2}\|x^{*}-x_0\|^{2},
\]
so that
\[
f(x_k)-f(x^{*})\le \frac{\sigma}{2\eta}\frac{\|x^{*}-x_0\|^{2}}{k}.
\]

\section*{Special Cases}
\begin{enumerate}
\item \textbf{Euclidean setup.} $h(x)=\tfrac12\|x\|^{2}$, $\sigma=1$, $D_{h}$ becomes $\tfrac12\|x-y\|^{2}$, and Theorem~\ref{thm:rate} coincides with the classical gradient descent rate.
\item \textbf{Negative entropy on the simplex.} $h(x)=\sum_{i}x_i\log x_i$, $\D{\,\cdot\,}{\,\cdot\,}$ is the Kullback–Leibler divergence; the algorithm reduces to exponentiated gradient descent with the same $O(1/k)$ guarantee.
\item \textbf{Adaptive stepsizes.} If $\eta_k = \frac{1}{L_k}$ with $L_k$ an estimate of the local relative smoothness constant, the same proof works by replacing $L$ with $L_k$ in Step 11, yielding a non‑increasing bound.
\end{enumerate}

\end{document}